\section{Conclusion}
\label{sec:conclusion}
Serving latency-sensitive functions from volunteer machines needs more than running tasks several times and keeping
sandboxes alive. Churn destroys the state that makes short functions fast, and untrusted owners make that state risky to
create. We argued that attesting a locked image reduces hostile owners to the stops and delays that churn already causes,
and that what remains is a scheduling problem: deciding which machines are attested, given a tenant's key and kept warm,
when all of that is slow to set up, lost at the end of every session, and a leak risk on every extra machine.
\voss{} plans this state every 30 seconds with spare capacity for lost machines, places each request by its own chance of
missing the deadline, judges machines by the work they return, and replaces failed schedulers through a provider-run
coordinator without an agreement protocol among volunteers. In the clustered architecture, schedulers on volunteer machines
are reached through a provider-run publish/subscribe broker; in the global architecture, the scheduler runs on provider
hardware.
In a calibrated simulator \voss{} keeps P99 under 3\,s in \ViabAPThreeVoss{} of \ViabACells{} and \ViabBPThreeVoss{} of \ViabBCells{} settings where simple baselines do not, lowers P99 by
\GreedyNoneMed--\GreedyNoneHigh\% against provisioning on demand, and keeps \VossSuccAlpha\% of requests successful when 20\% of owners drop work. The advantage is not
universal: on confidential virtual machines at high load the simple baseline is the only feasible system, and geographic
clusters, cost-aware caching and the reliability score did not pay for themselves. The next steps are the loss-model planner
of Section~\ref{sec:discussion} and a multi-machine testbed to validate the simulator and the networking design that this
paper models only as delay.
